% !TEX program = xelatex
\documentclass{英语试卷}
\begin{document}
\examtitle{2024 学年第一学期浙江省 9+1 高中联盟高二年级期中考试}{英语}
\par 英语试题\par
\par 命题：慈溪中学 宋浩颖 李晓燕 审题：桐乡高级中学 潘丽丽 义乌中学 王文伟\par
\parttitle{第一部分 听力（共两节，满分 30 分）}
\par 做题时，先将答案标在试卷上，录音内容结束后，你将有两分钟的时间将试卷上的答案转涂到答题卡上。\par
\sectitle{第一节（共5小题：每小题1.5分，满分7.5分）}
\par 听下面 5 段对话。每段对话后有一个小题，从题中所给的 A、B、C 三个选项中选出最佳选项。听完每段对话后，你都有 10 秒钟的时间来回答有关小题和阅读下一小题。每段对话仅读一遍。\par
\question{1. What does the woman recycle?}
\begin{choices}(3)
  \item Plastic.
  \item Glass.
  \item Paper.
\end{choices}
\question{2. What's wrong with Jane?}
\begin{choices}(3)
  \item She misses her home very much.
  \item She hasn't received her mother's letter.
  \item She is worried about her mother's health.
\end{choices}
\question{3. Where does the conversation take place?}
\begin{choices}(3)
  \item In a hospital.
  \item On a bus.
  \item In a classroom.
\end{choices}
\question{4. How will the speakers go to the bookstore?}
\begin{choices}(3)
  \item By taxi.
  \item By bus.
  \item By bike.
\end{choices}
\question{5. What does the woman say about the meeting?}
\begin{choices}(3)
  \item It has been brought forward.
  \item It has been called off.
  \item It will begin on time.
\end{choices}
\sectitle{第二节（共15小题：每小题1.5分，满分22.5分）}
\par 听下面 5 段对话或独白。每段对话或独白后有几个小题，从题中所给的 A、B、C 三个选项中选出最佳选项。听每段对话或独白前，你将有时间阅读各个小题，每小题 5 秒钟；听完后，各小题将给出 5 秒钟的作答时间。每段对话或独白读两遍。 听第 6 段材料，回答第 6、7 题。\par
\question{6. Why does the man like hiking?}
\begin{choices}(3)
  \item He can get close to nature.
  \item He can meet wild animals.
  \item He can make some friends.
\end{choices}
\question{7. How far will the man walk a day on average?}
\begin{choices}(3)
  \item 3 miles.
  \item 10 miles.
  \item 30 miles.
\end{choices}
\par 听第 7 段材料，回答第 8、9 题。\par
\question{8. What is the woman weak in?}
\begin{choices}(3)
  \item Reading in French.
  \item Writing in Spanish.
  \item Speaking in Spanish.
\end{choices}
\question{9. Why does the woman talk to the man?}
\begin{choices}(3)
  \item To apply for a job.
  \item To discuss her studies.
  \item To choose a language course.
\end{choices}
\par 听第 8 段材料，回答第 10 至 12 题。\par
\question{10. Why did the woman leave Los Angeles?}
\begin{choices}(3)
  \item To get a job.
  \item To go to college.
  \item To enrich her life.
\end{choices}
\question{11. What's the woman's next plan?}
\begin{choices}(3)
  \item Move to Chicago one day.
  \item Go to Asia in several months.
  \item Live in New York for a year.
\end{choices}
\question{12. Which city does the man want to settle in?}
\begin{choices}(3)
  \item Chicago.
  \item Los Angeles.
  \item New York City.
\end{choices}
\par 听第 9 段材料，回答第 13 至 16 题，\par
\question{13. How does the woman feel about rock music?}
\begin{choices}(3)
  \item It never gets boring.
  \item It comes from the heart.
  \item It's something you have to think about.
\end{choices}
\question{14. What did the woman look like?}
\begin{choices}(3)
  \item She had long black hair.
  \item She had black eyes.
  \item She wore leather clothes.
\end{choices}
\question{15. What was the name of the woman's band?}
\begin{choices}(3)
  \item Up Your Alley.
  \item The Blackhearts.
  \item Sly and the Family Stone.
\end{choices}
\question{16. What does the woman suggest in the end?}
\begin{choices}(3)
  \item She still wants to perform.
  \item People still don't understand her.
  \item She didn't actually have much influence.
\end{choices}
\par 听第 10 段材料，回答第 17 至 20 题。\par
\question{17. When did Jessica complete her world journey?}
\begin{choices}(3)
  \item In 1993.
  \item In 1996.
  \item In 2010.
\end{choices}
\question{18. What can we learn about Jessica Waston?}
\begin{choices}(3)
  \item She has written a book called True Spirit.
  \item She lost fingers while climbing a mountain.
  \item She traveled around the world alone on a boat.
\end{choices}
\question{19. What led Jordan to climb Qomolangma?}
\begin{choices}(3)
  \item A painting.
  \item Her parents.
  \item A book.
\end{choices}
\question{20. What is people's attitude to the parents of Jessica and Jordan?}
\begin{choices}(3)
  \item Supportive.
  \item Opposed.
  \item Unclear.
\end{choices}
\parttitle{第二部分 阅读理解（共两节，满分50分）}
\sectitle{第一节（共15小题：每小题2.5分，满分37.5分）}
\par 阅读下列短文，从每题所给的 A、B、C 和 D 四个选项中，选出最佳答案。\par
\passagelabel{A}
\par Including winners of the Pulitzer Prize, Edgar Award, and more, these books have been read, judged, and celebrated by the best this year! Pick up an award-winning story for any age.\par
\par Vinner of the Edgar Award\par
\par Vera Wong's Unsolicited Advice for Murderers(杀人犯) Jesse Q. Sutanto\par
\par Paperback \$18.00 Add to cart\par
\par Vera Wong, a lonely old lady who lives above her forgotten tea shop in the middle of San Francisco's Chinatown, walked downstairs to find a curious thing—a dead man in the middle of her tea shop. This lonely shopkeeper takes it upon herself to solve a murder in the most unusual way in this best-selling novel.\par
\par\vspace{0.4em}
\begin{center}\begin{tabular}{{|l|p{\dimexpr\linewidth-3em-4\tabcolsep}|}}
\hline
Winner of the Pulitzer in NonfictionHis Name Is George FloydRobert Samuels \&amp; Toluse OlorunnipaPaperback \$20.00 Add to cart & A landmark biography that reveals how racism(种族主义) shaped George Floyd\&\#x27;s life—from his family\&\#x27;s roots in the tobacco fields, to ongoing inequality in housing, education, and health care—telling the story of how one man\&\#x27;s tragic experience brought about a global movement for change. \\
James Beard Award WinnerSavor(品味)Fatima AliHardcover \$28.00 Add o cart & A young cook whose dreams were cut short savors every last minute as she explores food and adventure, illness and death in Savor, an inspiring autobiography and family story that sweeps from Pakistan to Manhattan and beyond. At its core, it is a story about what it means to truly live. \\
\hline
\end{tabular}\end{center}
\vspace{0.4em}
\question{21. Who would probably be interested in Vera Wong's Unsolicited Advice for Murderers?}
\begin{choices}(2)
  \item Mystery fans.
  \item Tea shop owners.
  \item Professional cooks.
  \item Chinatown visitors.
\end{choices}
\question{22. What do three books share in common?}
\begin{choices}(2)
  \item They're now at a discount.
  \item They explore the same theme.
  \item They're award winners of this year.
  \item They are based on true life stories.
\end{choices}
\question{23. Where can you find the text?}
\begin{choices}(2)
  \item A website.
  \item A magazine.
  \item A textbook.
  \item A newspaper.
\end{choices}
\passagelabel{B}
\par "Why is every superhero movie an origin story?" complained Entertainment Weekly film critic Adam Markovitz after seeing a trailer(预告片) for this summer’s Man of Steel. Perhaps we love origin stories, Markovitz suggested, because they "show the exact moment when a normal guy goes from being just like us to being somehow better, faster, stronger."\par
\par I tend to disagree. I think origin stories show us not how to become super but how to be heroes, choosing other people's needs over the pursuit of one's own wealth and power. I've learned this through hundreds of conversations at comic book meetings, where fans have been remarkably honest about their lives and the inspiration they draw from superhero stories.\par
\par I've found that superheroes go through three types of life-changing experiences that we can relate to. The first is trauma(创伤), which lies at the heart of Batman's origin story, in which Bruce Wayne commits himself to fighting crime after seeing his parents killed. In real life, many people experience "stress-caused growth" after a trauma and determine to help others.\par
\par The second life-changing force is destiny(命运). Consider Buffy the Vampire Slayer, about a normal teenager who discovers she's the "Chosen One"—born with supernatural powers to fight demons. Buffy is unwilling to accept her destiny, yet she throws herself into her new job.\par
\par Lastly, there's pure chance, which transformed a young Spider-Man, who was using his power for selfish(自私的) purposes until his beloved uncle was killed by a street thug. Spider-Man's heroism is an example of how random negative events cause many of us to reflect on our lives and choose a different path.\par
\par At their best, superhero origin stories inspire us and provide models of coping with difficulties, finding meaning in loss and trauma, discovering our strengths and using them for good purpose. (Wearing a cape or tights is not a must.)\par
\question{24. How did Adam Markovitz feel about every superhero movie being an origin story?}
\begin{choices}(2)
  \item Surprised.
  \item Disappointed.
  \item Fascinated.
  \item Inspired.
\end{choices}
\question{25. What experience led Bruce Wayne to become a hero?}
\begin{choices}(2)
  \item Suffering from great stress.
  \item Fighting crime with great devotion.
  \item Relating to Batman's origin story.
  \item Witnessing his parents being killed.
\end{choices}
\question{26. What did Buffy do after she found her destiny?}
\begin{choices}(2)
  \item She couldn't wait to accept her role.
  \item She used her power for selfish purposes.
  \item She refused to take on the responsibility.
  \item She committed herself to her role eventually
\end{choices}
\question{27. According to the author, why do people love origin stories?}
\begin{choices}(2)
  \item Because they inspire us to put others before ourselves.
  \item Because they lie at the heart of every superhero movie.
  \item Because they explore the impact of negative events on life.
  \item Because they show how ordinary people become extraordinary.
\end{choices}
\passagelabel{C}
\par In the park, there are 9485 of them. You sit on them. To rest. Read a book. Wait for a friend. Or it's such a day when you have nothing to do, and this is a delicious place to accomplish absolutely nothing. Or you can drift off and muse on the plaque(纪念牌) attached there, representing a story behind the bench. The Central Park bench You aren't just sitting on wood. You are sitting on memories.\par
\par Caitlin LaMorte was the development manager for the Adopt-a-Bench program which began in 1986 as a way to finance the maintenance of the benches. If you can afford it, it's simple enough. Pay \$10000 and you get to put a plaque on a bench, saying almost whatever you want, up to a suggested maximum of four lines of 30 characters each. And then it's there forever. Up to now,4223 of the benches have been adopted. While plenty of benches remain unadopted, some areas are sold out. For instance, all the benches facing the lake.\par
\par Benches can serve many purposes. Quite often, they're adopted to remember a relative or friend. Or on occasion, a pet. Increasingly, Ms. LaMorte said, "we have more plaques that are happy." Graduations or birthdays or birth wishes or wedding gifts.\par
\par Last year, Karen May figured she would adopt a bench for her children. When her elder son, Theodore, proposed to his girlfriend, Lucinda, he wrote a brief children's book for her called "Two Red Foxes," because foxes were a recurring theme in her upbringing (as when a fox sped across the field when her father proposed to her mother). The foxes were stand-ins for Ms. May's son and his bride-to-be. And there was the fact that her younger son, Thornwell, went by the nickname Pup. So there is a bench facing the lake that reads, "Two Red Foxes and a Pup." People will sit there and not know what it means. But Ms. May knows.\par
\question{28. What is the Adopt-a-Bench program intended for?}
\begin{choices}(2)
  \item Raising funds for bench care.
  \item Placing more plaques on benches.
  \item Marketing the unadopted benches.
  \item Attracting visitors to the central park.
\end{choices}
\question{29. What does the underlined word "stand-ins" probably mean?}
\begin{choices}(2)
  \item Relatives.
  \item Wishes.
  \item Symbols.
  \item Gifts.
\end{choices}
\question{30. Why did Karen May adopt the bench?}
\begin{choices}(2)
  \item To remember two red foxes and a puppy.
  \item To secure a lakeside bench for family use.
  \item To promote her son's newly published book.
  \item To celebrate her children's special moments.
\end{choices}
\question{31. Which of the following can be the best title for the text?}
\begin{choices}(2)
  \item Adopt-a-Bench Program
  \item Two Red Foxes and a Pup
  \item Benches With Stories to Tell
  \item A Creative Gift in the Park
\end{choices}
\passagelabel{D}
\par Two small white dogs appeared as I arrived at a ranch house a few miles off the highway in Hempstead, followed by John Mendola, a retired police officer. He introduced me to the dogs, Princess Ariel and Princess Jasmine. They were named for a deceased dog named Princess to whom they are similar in every detail. And they should be: they are Princess's clones.\par
\par It has now been nearly thirty years since cloning mammals(哺乳动物) became possible. The technology has mainly been used to produce cattle, sheep, and pigs. However, since 2005, more than two thousand dogs have been successfully cloned. Biologically, it is not very different from that of cloned cows or sheep, but in other respects the cloning of pets raises a host of ethical(伦理的) and identity issues\par
\par There are issues of exploitation(剥削)—both of grieving owners who hope to bring back their lost pets, and of the unseen animals whose bodies are used in making a clone. There's the issue of supply: the production of tailored dogs in a society when so many good, naturally born ones in shelters are in need of adoption. Finally, there's an existential issue: who, exactly, is produced when a dog is cloned? Dogs are seen by most owners as a species of individuals, with distinct personalities. We keep them for companionship. As a result, we have placed our ideas of selfhood onto them, giving them preferences, fears, and moods. But, if it is dogs' individuality that we value, how can we accept the idea that their unique and unreproducible selves can, in fact, be reproduced?\par
\par As a matter of fact, among the people who have cloned their dogs, many struggled to say exactly what it was about the original animal that they had wanted to reproduce, especially in contrast to other dogs they had loved but hadn't felt like cloning. Many spoke of the original as simply "special" —but the specific nature of that specialness seemed almost impossible to be explained. Perhaps the desire was less about re-creating the original dog than about bringing back the unique relationship that had been established with the animal.\par
\question{32. What can we learn about the two white dogs in the first paragraph?}
\begin{choices}(2)
  \item They were named after a princess.
  \item They were clones of a late dog.
  \item They were distinct from each other.
  \item They were naturally produced by Princess.
\end{choices}
\question{33. How is the cloning of dogs different from that of other mammals?}
\begin{choices}(2)
  \item It has higher success rates.
  \item It raises more moral concerns.
  \item It is technically more challenging.
  \item It stirs greater public curiosity.
\end{choices}
\question{34. Which word best describes the author's attitude to the cloning of dogs?}
\begin{choices}(2)
  \item Unclear.
  \item Objective.
  \item Favourable.
  \item Disapproving.
\end{choices}
\question{35. Why do people desire to clone their dogs?}
\begin{choices}(2)
  \item To deal with the identity issue.
  \item To reproduce the original animal.
  \item To recreate a distinctive emotional bond
  \item To find out the specialness of the original dog.
\end{choices}
\sectitle{第二节（共5小题；每小题2.5分，满分12.5分）}
\par 阅读下面短文，从短文后的选项中选出能填入空白处的最佳选项。选项中有两项为多余选项，\par
\par Everybody lies to some degree. \blanknum{36} However, nobody wants to be caught. Of course there's a huge gap that separates little white lies from the whoppers(弥天大谎), but learning how to tell if someone is lying is a skill that'll always come in handy. Here's what to watch for.\par
\par \blanknum{37}\par
\par There are many signs someone could be lying. For example, one may make too little or too much eye contact, sweat or flush in their face, and have difficulty controlling the volume and tone of their voice. \blanknum{38} Hand gestures that normally accompany talking may occur with less frequency. Arm and leg movements may also be fewer. It seems that liars would become more concentrated on telling the lie, so they get quieter in their body.\par
\par Ask questions——quickly.\par
\par If you suspect you're being lied, try this technique, which experts say can trip up a liar. Try asking questions quickly—one after the other. "The first lie is easy," explains Kang Lee, professor at the Ontario Institute for Studies in Education's Department of Applied Psychology and Human Development. "\blanknum{39} When you continue to ask questions and put people on the spot, it gets harder to maintain the lie."\par
\par Check for emotional "leaks".\par
\par Micro-expressions often show a person's real thoughts. If you were to watch people very closely on videotape, you would see them showing their true emotion just before they show the fake expression designed to cover up the lie. But these extremely brief facial movements, some only lasting a quarter of a second, aren't easy to spot. \blanknum{40} And intentional liars tend to add other expressions, like smiling, to hide a lie.\par
\par So here's hoping that the next time someone throws a lie your way, you'll know just how to catch it.\par
\begin{choices}(1)
  \item Listen to the voices.
  \item Get better at body language
  \item The follow-up lie is more difficult.
  \item Interestingly, people who are lying often become more still.
  \item Even profes ionals trained in lie detection can't always tell them apart.
  \item In some cultures, it's considered inappropriate to maintain eye contact.
  \item Lies occur between friends, teachers and students, husbands and wives.
\end{choices}
\parttitle{第三部分 语言知识运用（共两节，满分 30 分）}
\sectitle{第一节（共15小题；每小题1分，满分15分）}
\par 阅读下面短文，从短文后各题所给的 A、B、C 和 D 四个选项中，选出可以填入空白处的最佳选项，\par
\par "You’ve lost 60\% of your hearing in your right ear and it's permanent(永久的)." "What?!" I stared at the doctor. So he \blanknum{41} himself, louder this time. He thought I couldn't \blanknum{42} him, but I had. I was just in \blanknum{43}. How had this happened? How am I supposed to live my life with one ear? I burst out \blanknum{44}.\par
\par Right before Thanksgiving 2023, I got what I thought was a \blanknum{45} head cold. No pain, No cough. It just made my right ear feel as if it was underwater. I told people to "talk to my \blanknum{46} ear." I figured it would resolve on its own \blanknum{47}.\par
\par It didn't. One month went by, then two, and I \blanknum{48} couldn't hear well, but I didn't get it \blanknum{49}. It wasn't until the end of January, when I went to an ear, nose and throat (ENT) doctor for an unrelated issue, that I thought to \blanknum{50} it. As soon as I did, the doctor looked at me in alarm and told me to make a (an) \blanknum{51} for a hearing test the next day. After much testing, I was \blanknum{52} as having SSHL, a medical term for "I suddenly lost my hearing for no apparent reason."\par
\par A month later, I was fitted with my new hearing \blanknum{53}. I hadn't realized how much I was \blanknum{54} with my damaged hearing until the moment it was given back to me. I burst into happy tears. I was so grateful for all the medical \blanknum{55} and technology that gave me back my hearing.\par
\cz{41}{introduced}{excused}{bent}{repeated}
\cz{42}{hear}{see}{understand}{accept}
\cz{43}{amazement}{trouble}{shock}{fear}
\cz{44}{crying}{laughing}{running}{cheering}
\cz{45}{critical}{minor}{awesome}{painful}
\cz{46}{right}{nice}{good}{poor}
\cz{47}{permanently}{fortunately}{eventually}{hopefully}
\cz{48}{even}{still}{yet}{already}
\cz{49}{checked out}{lifted up}{looked into}{brought about}
\cz{50}{explain}{warn}{investigate}{mention}
\cz{51}{break}{appointment}{preparation}{survey}
\cz{52}{interpreted}{labelled}{predicted}{identified}
\cz{53}{ability}{aid}{relief}{wear}
\cz{54}{dealing}{struggling}{missing}{enjoying}
\cz{55}{professionals}{insurance}{sensors}{knowledge}
\sectitle{第二节（共10小题；每小题1.5分，满分15分）}
\par 阅读下面短文，在空白处填入 1 个适当的单词或括号内单词的正确形式，\par
\par StarLink, developed by SpaceX, intends to provide high-speed internet access globally. This satellite-based service \blanknum{56} (make) up of a series of low Earth orbit(LEO) satellites \blanknum{57} communicate through advanced laser technology, forming a global network.\par
\par StarLink is designed to be a solution for developing regions, offering \blanknum{58} (rely) internet resistant to weather conditions, with speeds \blanknum{59} (expect) to reach up to 1 Gbps—significantly faster than traditional broadband. Initially \blanknum{60} (focus) on the U.S. and Canada, the service plans to expand worldwide by the end of 2022, aiming \blanknum{61} (bridge) the digital divide between urban and rural areas.\par
\par The project not only seeks to enhance connectivity for those in challenging circumstances but also promises to provide more people with access \blanknum{62} information and resources. \blanknum{63} it progresses, StarLink is ready to transform the landscape of internet access, \blanknum{64} (potential) changing how people around the world connect and interact. In summary, StarLink represents \blanknum{65} groundbreaking advancement in global internet accessibility, with significant implications for individuals and communities traditionally left behind.\par
\parttitle{第四部分 写作（共两节，满分 40 分）}
\sectitle{第一节（满分15分）}
\par 假定你是李华，本月新西兰学生代表团正在访问你校，你将作为学生代表介绍学校的社团活动。请为此写一篇发言稿，内容包括：\par
\par (1) 表示欢迎社团; (2) 活动介绍; (3) 推荐体验的社团。\par
\par 注意：(1) 写作词数应为 80 左右；(2) 请按如下格式在答题纸的相应位置作答。\par
\par\vspace{0.4em}
\begin{center}\begin{tabular}{{|l|}}
\hline
Dear friends from New Zealand, \\
\_ \\
\_ \\
\hline
\end{tabular}\end{center}
\vspace{0.4em}
\sectitle{第二节（满分25分）}
\par 阅读下面材料，根据其内容和所给段落开头语续写两段，使之构成一篇完整的短文。\par
\par "The cleaners said I could pick our stuff today," Jet's mom said on the drive to school.\par
\par "Do you think they were able to clean Big Dog?" Jet asked, trying not to sound too worried.\par
\par But he was worried. Big Dog was his favourite stuffed animal. Grandpa had won it for him at the fair when he was five years old. Jet was eight now, and he didn't need to sleep with it or anything. But he did need it. He needed it to sit on the shelf above his bed—the shelf that was now charred(烧焦的) by the house fire.\par
\par "I'm not sure what they were able to save," Mom said. "They said they weren't able to save our photo albums. I sure will miss those pictures." Jet could tell she was really sad.\par
\par "I'll pick you up after school, and we'll go through the boxes together. Then we'll see if your Big Dog can come back to us." Mom said before she hugged him good-bye and drove off.\par
\par Jet tried to focus on his schoolwork all day, but he couldn't keep his mind off Big Dog. He clutched(紧抓) his backpack through the halls as if he were carrying his old stuffed friend. During library time, he checked out a book about dogs. At break time, his friends were playing games while his mind was completely occupied by Big Dog.\par
\par RRRRRRINNNGGG! Time for art class.\par
\par Mr. Light said it was "artist's choice" day, so they could work on whatever they wished.\par
\par Jet stared at an empty page. He thought of his mom picking up the boxes of things that had been damaged by the fire. Then he remembered how sad his mom was about the photo albums. Most things could be replaced at the store. But the pictures were gone forever. If only the photo album could somehow be saved, Jet thought. He wished there was a way to get those memories back, to bring a smile to his mom's face again.\par
\par 注意:\par
\par (1) 续写词数应为 150 左右;\par
\par (2) 请按如下格式在答题卡的相应位置作答。\par
\writingpara{Paragraph 1: Then, all of a sudden, he knew what he could draw.}
\writinglines{8}
\writingpara{Paragraph 2: Jet gave his mom a huge hug, and said, "I have a surprise for you too!".}
\writinglines{8}
\end{document}
